\documentclass[main]{subfiles}
\usepackage[sols]{workbook}
\numbering{ws1-2}

\begin{document}

\ifSubfilesClassLoaded{%
  \chaptermark{\chapterone}
}{}

\section{Density and the Definite Integral} \label{ws1-2}

\begin{framed}
\textbf{Lesson Goals:}
\begin{itemize}
\item You should be able to apply the \textit{Slice} $\rightarrow$ \textit{Approximate} $\rightarrow$ \textit{Sum} $\rightarrow$ \textit{Limit} strategy to find the total mass/amount over a 2D region with a given density function.
\item You should be able to decide how to draw, label, orient, and slice a 2D region with a given density function, and apply various geometric techniques to approximate the amount within each slice.
\item You should understand Riemann sum notation, and the definition of the definite integral as a limit of a Riemann sum.
\end{itemize}


\end{framed}


\begin{enumerate}
\item
 \label{spring-2015-midterm1-2d-density-slicing} 
A sail boat has a main sail that is a right triangle.  After a windy season of sailing in the San Fransisco Bay, the sail has been splashed over and over by the bay water, and the salt in the water has dried into the sail.
 

\parbox[t]{4.5in}{
The bottom edge of the sail is $2$m long, and the height of the sail is $6$m.  The density of salt in the fabric of the sail is given by $\rho(x)$ grams per square meter, where $x$ is the distance from the bottom edge of the sail in meters. \\
      
      \begin{enumerate}
      \item Write a Riemann sum that gives the approximate total mass of salt that is in the fabric of the sail. Please define any notation you use.
      \end{enumerate}
}  \quad
\parbox[t]{2in}{\vspace{-0.15in}
\includegraphics[width=2in]{ws/images/spring2015-sail .jpg}}
 
\pspace{\vfill}
\begin{solution}

\parbox[t]{2in}{\vspace*{0in} \includegraphics[width=1.5in]{ws/images/spring2015-sail-slicing .png}}
\parbox[t]{4.5in}{
We slice the interval from $x_0=0$ at the bottom of the sail to $x_n=6$ into $n$ slices of equal width, $\Delta x = \frac{6}{n}$, so that each slice has approximately uniform density.  The $k^{th}$ slice is at a height $x_k = 0 + k\Delta x$ meters from the bottom of the sail, and has a width $w_k$ as drawn in the figure.

By comparing similar triangles (the small one above the $k^{th}$ slice, and the whole sail), we can see that $\displaystyle \frac{w_k}{6-x_k} = \frac{2}{6} = \frac{1}{3}$, and so $\displaystyle w_k = \frac{6-x_k}{3} = 2-\frac{x_k}{3}$.

The mass of salt on the $k^{th}$ slice is given approximately by

\begin{eqnarray*}
\mathrm{mass} &\approx& \mathrm{density} \cdot \mathrm{volume}\\
&\approx& \rho (x_k) \cdot w_k \cdot \Delta x \\
&=& \rho(x_k)\left(2-\frac{x_k}{3}\right) \Delta x
\end{eqnarray*}

The total mass of salt in the sail is therefore given by the sum over all $n$ slices, hence:

\[ \boxed{\mathrm{total\; mass} \approx \sum_{k=1}^n \rho(x_k)\left(2-\frac{x_k}{3}\right) \Delta x }\]
}
\end{solution}


\begin{enumerate}
\item[(b)]
\begin{problem}[1.25in]
By taking an appropriate limit, write a definite integral that gives
      the exact mass of salt in the sail.  Do not evaluate this integral, but please write out explicitly the limit you are taking.
\end{problem}

\begin{solution}
The exact mass of the salt is the limit of the Riemann sum:
\[\mathrm{total\; mass} = \lim_{n\rightarrow \infty} \left(  \sum_{k=1}^n \rho(x_k)\left(2-\frac{x_k}{3}\right) \Delta x \right) = \boxed{\int_0^6 \rho(x) \left(2-\frac{x}{3}\right)\,dx} \]
\end{solution}

\end{enumerate}


\probonly{\newpage}

\begin{framed}
\centerline{\textbf{The Limit Definition of the Definite Integral}}

$$\boxed{\int_a^b f(x)\ dx = \lim_{n \to \infty}\Bigg[ \sum_{k=1}^n f(x_k)\Delta x\Bigg]}$$
where $\boxed{\Delta x = \frac{b-a}{n}}$ and $\boxed{x_k=a+k\Delta x}$. In this notation, $[a, b]$ represents the interval of $x$-values that is being sliced, $n$ represents the number of slices we're using, $\Delta x$ represents the width/thickness of each slice, and $x_k$ represents the location of the $k$-th slice.

\end{framed}
  
\item  \label{density-concentric-rings-1}

  People in the Boston area like to live near the city center, so the
  population density around Boston is $\rho(r) = \frac{36,000}{r^2 + 2r +
    1}$ people per square mile, where $r$ is the distance in miles to the
  center of Boston.  We'd like to find the number of people who live
  within 5 miles of the center of Boston.

  \begin{enumerate}
  \item
    \begin{problem}[\vfill]
      Show in a sketch how you would \textbf{slice} the region in question.
    \end{problem}    

    \begin{solution}
      The region we are talking about (within 5 miles of the center of
      Boston) is the area enclosed by a circle, with the highest
      population density at the center.  Here is a plot showing the
      population density in the city (red represents the highest density,
      and purple represents the lowest):
      \begin{center}
        \includegraphics[height=1.25in]{ws/images/density-concentric-rings-1-1}
      \end{center}
      As always, we want to slice the region so that the population
      density in each slice is almost constant.  In this case, we can
      accomplish that by slicing into concentric rings.  (All of the
      points within a single ring are approximately the same distance from
      the city center, so they have approximately the same population
      density.)  This means that we are slicing the interval $[0, 5]$ into
      $n$ pieces of equal width, and we'll label the endpoints $r_0 = 0,
      \ldots, r_n = 5$.
      \begin{center}
        \includegraphics[height=1.5in]{ws/images/density-concentric-rings-1-2}
      \end{center}
    \end{solution}

  \item
    \begin{problem}[\vfill] \label{density-concentric-rings-1-kth-slice} How
      would you \textbf{approximate} the number of people who live in the
      $k$-th slice?
    \end{problem}    

    \begin{solution}
      The $k$-th slice is a ring with outer radius $r_k$ and inner radius
      $r_{k-1}$:
      \begin{center}
        \includegraphics[height=1in]{ws/images/density-concentric-rings-1-slice}
      \end{center}
      The population density in the $k$-th slice is approximately equal to
      $\rho(r_k)$.  We can think of the $k$-th slice roughly as a rectangle that has been curled into a ring. If we ``unrolled" this ring into a rectangle, the circumference of the outer edge of the ring (which is $2 \pi r_k$) would give the approximate length of the rectangle, and the thickness of the ring (which is $\Delta r$) would give the approximate width of the rectangle. So, the area of this slice is approximately $2 \pi r_k\Delta r$.  
      
      Therefore, the number of people living in the $k$-th slice is approximately $\rho(r_k) \cdot 2 \pi r_k \Delta r = \boxed{2 \pi r_k
        \rho(r_k) \Delta r}$.
        
 Notice that we could have found the \textit{exact} area of the $k$-th slice by subtracting the area of the two disks that bound it, which would give $\pi(r_k)^2-\pi(r_{k-1})^2$=$\pi\big[(r_k)^2-(r_{k-1})^2\big]$. However, this expression does not explicitly contain $\Delta r$, which would make writing a Riemann sum in the next step much more difficult. So here, it is better to \textit{approximate} the area of the $k$-th slice with a rectangle than to find it exactly. Also, as our slices get thinner, the $k$-th ring will be better approximated by a rectangle, since the circumferences of the inner and outer circles will be closer in length, better resembling the equal parallel sides of a rectangle.
        
        
    \end{solution}

  \item
    \begin{problem}[\vfill]
      Write a general Riemann \textbf{sum} that approximates the number of
      people who live within 5 miles of the center of Boston.
    \end{problem}    

    \begin{solution}
      Summing our answer from \ref{density-concentric-rings-1-kth-slice}
      over all slices, the number of people living within 5 miles of the
      center of Boston is approximately \fbox{$\displaystyle \sum_{k=1}^n
        2 \pi r_k \rho(r_k) \Delta r$}.
    \end{solution}

  \item
    \begin{problem}[\vfill]
      By \textbf{taking an appropriate limit}, find a definite integral
      that gives the number of people who live within 5 miles of the
      center of Boston.
    \end{problem}

    \begin{solution}
      Taking the limit as $n \to \infty$ of the general Riemann sum we
      found gives the integral $\boxed{\int_0^5 2 \pi r \rho(r)~dr}$. 
    \end{solution}
  \end{enumerate}
  
  
  
  
%  \probonly{\newpage}
  
  
  
  
\item 
\begin{problem}[10in]
Evaluate $\displaystyle \lim_{n \to \infty} \sum_{k=1}^n \frac{1}{5+k\Delta x}\, \Delta x$, where $\Delta x = \frac{4}{n}$
by interpreting it as a definite integral and evaluating the integral.

(There are two such integrals you can come up with. Both, of course, have the same value.)
\end{problem}

\begin{solution}
As the problem mentions, there are two possible definite integrals that could be defined by the limit of the Riemann sum given above, depending on what you interpret to be $x_k$ in the summation. By relating $\displaystyle  \sum_{k=1}^n \frac{1}{5+k\Delta x}\, \Delta x$ to $\displaystyle  \sum_{k=1}^n f(x_k)\, \Delta x$ in the limit definition of the integral, we can see that $$f(x_k)= \frac{1}{5+k\Delta x}$$ 

However, there are two possible choices for what $x_k$ could then be. It could be that $x_k=5+k\Delta x$; this would imply that $ f(x_k)=\frac{1}{x_k}$, and thus $ f(x)=\frac1x$. Or it could be that $x_k=k\Delta x$; this would imply that $f(x_k)= \frac{1}{5+x_k}$, and thus $f(x)= \frac{1}{5+x}$. We'll explore both options here:

\uline{\textbf{Option 1:}} $\displaystyle f(x)=\frac{1}{x}$, $x_k=5+k\Delta x$, and $\Delta x=\displaystyle\frac4n$.

We can always find the endpoints of our slicing interval by taking $x_k$ and setting $k=0$ and $k=n$, since $x_0$ gives the left endpoint of the first slice, and $x_n$ gives the right endpoint of the last slice. Thus, 
\begin{align*}
a&=x_0=5+0\cdot\Delta x=5\\
b&=x_n=5+n\Delta x=5+n\cdot\frac{4}{n}=5+4=9
\end{align*}

Therefore,
$$\displaystyle \lim_{n \to \infty} \sum_{k=1}^n \frac{1}{5+k\Delta x}\, \Delta x = \lim_{n \to \infty} \sum_{k=1}^n \frac{1}{x_k}\, \Delta x = \int_5^9 \frac1x\ dx$$

Evaluating this definite integral, we get:

$$\int_5^9 \frac1x\ dx = \ln|x|\Bigg|_{x=5}^{x=9} = \ln(9)-\ln(5)=\boxed{\ln\Bigg(\frac95\Bigg)}$$

\uline{\textbf{Option 2:}} $\displaystyle f(x)=\frac{1}{5+x}$, $x_k=k\Delta x$, and $\Delta x=\displaystyle\frac4n$.

Finding the endpoints of the interval using the same method as before, we get:
\begin{align*}
a&=x_0=0\cdot\Delta x=0\\
b&=x_n=n\Delta x=n\cdot\frac{4}{n}=4
\end{align*}

Therefore,
$$\displaystyle \lim_{n \to \infty} \sum_{k=1}^n \frac{1}{5+k\Delta x}\, \Delta x = \lim_{n \to \infty} \sum_{k=1}^n \frac{1}{5+x_k}\, \Delta x = \int_0^4 \frac{1}{5+x}\ dx$$

We can evaluate this integral using the substitution $u=5+x$ and $du=dx$. So, when $x=0$, $u=5$; and when $x=4$, $u=9$. This gives us the equivalent integral:
$$\int_5^9 \frac{1}{u}\ du$$
However, this integral is exactly the same as the one we found in Option 1; the only difference is that the dummy variable is now $u$ instead of $x$. Thus, we get the same final answer, $\boxed{\ln\Bigg(\frac95\Bigg)}$.
\end{solution}





 % \probonly{\newpage}
  
  
  
  
  
  
\item  \label{density-triangle-vegetation}

  A plot of land is shaped like an equilateral triangle with sides of
  length 2 miles.  One side of the plot lies along a river, and vegetation
  density varies with $x$, distance to the river.  Suppose vegetation
  density is given by $\rho(x)$.

  \begin{enumerate}
  \item
    \begin{problem}[5.5in]
      Write a general Riemann sum that approximates the amount of
      vegetation in the plot.
    \end{problem}

    \begin{solution}
      Here is a sketch (the blue line at the bottom represents the
      river):
      \begin{center}
        \includegraphics[width=1.5in]{ws/images/density-triangle-vegetation}
      \end{center}
      Notice that, since $x$ represents distance to the river, the
      ``$x$-axis'' in this picture actually extends up.  Since the
      equilateral triangle has sides of length 2 miles, its height is
      $\sqrt{3}$ miles.  (If you're not convinced, draw the height
      straight down the middle of the triangle and then use the
      Pythagorean Theorem on the resulting right triangles.)

      In order to have slices in which the vegetation density is
      approximately constant, we should slice parallel to the river.
      Let's make $n$ slices of equal thickness $\Delta x$ (slicing the
      interval $[0, \sqrt{3}]$), label the $x$-values where we've sliced
      $x_0 = 0, \ldots, x_n = \sqrt{3}$, and focus on the $k$-th slice:
      \begin{center}
        \includegraphics[width=1.5in]{ws/images/density-triangle-vegetation-sliced}
      \end{center}
      The vegetation density throughout this slice is approximately
      $\rho(x_k)$, so we need to approximate the area of this slice.  The
      slice is approximately a rectangle of height $\Delta x$.  Let's call
      its width (at the top) $w_k$, so that the area of the $k$-th slice
      is approximately $w_k \Delta x$.
      \begin{center}
        \includegraphics[width=1.5in]{ws/images/density-triangle-vegetation-kth-slice}
      \end{center}
      Since we are slicing $x$, we need to write $w_k$ in terms of $x_k$
      to get a meaningful Riemann sum.  We can do this using similar
      triangles.  We know each side of our original triangle has length 2,
      and the height of the original triangle is $\sqrt{3}$.

      If we focus on the smaller equilateral triangle (the one whose base
      is $w_k$), we see that its width is $w_k$ and its height is
      $\sqrt{3} - x_k$ (remember that $x$ measures from the bottom of the
      triangle in our picture).  So, using similar triangles,
      $\displaystyle \frac{\sqrt{3}}{2} = \frac{\sqrt{3} - x_k}{w_k}$.
      Therefore, $\displaystyle w_k = \frac{2 (\sqrt{3} -
        x_k)}{\sqrt{3}}$.

      So, the area of the $k$-th slice is approximately $\displaystyle w_k
      \Delta x = \frac{2 (\sqrt{3} - x_k)}{\sqrt{3}} \Delta x$.  Since the
      vegetation density in this slice is approximately $\rho(x_k)$, the
      amount of vegetation in this slice is approximately $\displaystyle
      \frac{2 (\sqrt{3} - x_k)}{\sqrt{3}} \rho(x_k) \Delta x$.  Summing
      over all slices, we get the Riemann sum approximation
      \fbox{$\displaystyle \sum_{k=1}^n \frac{2 (\sqrt{3} -
          x_k)}{\sqrt{3}} \rho(x_k) \Delta x$} for the amount of
      vegetation in the plot.
    \end{solution}

  \item
    \begin{problem}[1in]
      Find a definite integral that gives the exact amount of vegetation
      in the plot.
    \end{problem}

    \begin{solution}
      We simply take the limit as $n \to \infty$ of our Riemann sum
      approximation and remember that we had sliced the interval $[0,
      \sqrt{3}]$ originally.  This gives the integral \fbox{$\displaystyle
        \int_0^{\sqrt{3}} \frac{2 (\sqrt{3} - x)}{\sqrt{3}} \rho(x)~dx$}.
    \end{solution}
  \end{enumerate}
\pspace{\newpage}
  \item \begin{problem}[\vfill]
    A circular state park has radius 5 miles, as shown in the picture
    below.  A stream of length 8 miles runs through the park (the dotted
    line in the picture).  
    \begin{center}
      \includegraphics[width=1in]{ws/images/density-circle-chord-1}
    \end{center}
    The area northwest of the stream is
    mountainous, while the area southeast of the stream is forest.  (The
    forest is the shaded area in the picture below.)
    \begin{center}
      \includegraphics[width=1in]{ws/images/density-circle-chord-2-forest}
    \end{center}    
    The density of trees in the forest varies with distance from the
    stream.  Let's say the density is given by $\rho(s)$ trees per square
    mile, where $s$ measures distance (in miles) from the stream.  Write
    an integral that gives the number of trees in the forest.
    
   \probonly{\textit{Note:} Whatever variable you decide to slice, think
    carefully about its meaning and how it relates to the density
    function!} 
  
  \end{problem}



  \begin{solution}
    As usual, we want to slice so that the density in each slice is
    approximately constant.  Since density varies with distance to the
    stream, this means we need to slice parallel to the stream.  Let's
    rotate the picture so that the stream is vertical and put the entire
    picture on the $xy$-plane:
    \begin{center}
      \includegraphics[width=1.25in]{ws/images/density-circle-chord-2-coord}
    \end{center}
    Here, we've chosen to place the circle so that its center is at the
    origin.\footnote{This is a totally arbitrary choice, just as we chose
      arbitrarily to rotate the picture to make the stream vertical.}
    Having made this choice, we can describe the circle by the equation
    $x^2 + y^2 = 25$.  Let's figure out the coordinates of the stream.
    Since the stream is 8 miles long, we know that the top point of the
    stream has $y$-coordinate equal to 4, and the bottom point has
    $y$-coordinate equal to $-4$.  Since both points are on the circle
    $x^2 + y^2 = 25$, their $x$-coordinates satisfy $x^2 + 16 = 25$, or
    $x^2 = 9$.  Since we've drawn the stream on the left of the origin, we
    want the negative square root $x = -3$.  So, the stream is the line $x
    = -3$, the bottom point on the stream is $(-3, -4)$, and the top point
    is $(-3, 4)$.

    As we explained already, we should slice parallel to the stream; with
    the picture we've drawn, that means we should slice vertically.  Let's
    do the integral in terms of $x$\footnote{This is also an arbitrary
      choice.  You could certainly slice $s$ instead.} (which is
    \textbf{not} the same as $s$: $s$ measures distance to the stream,
    while $x$ is just a variable we've made up to label our picture).  So,
    we slice the interval $[-3, 5]$ into $n$ pieces of equal width $\Delta
    x = \frac{8}{n}$, label $x_0 = -3, \ldots, x_n = 5$ as usual (so $x_k
    = -3 + k \Delta x$), and focus on the $k$-th slice:
    \begin{center}
      \includegraphics[width=1.75in]{ws/images/density-circle-chord-1-right-slice}
    \end{center}
    We were told in the problem that the density of trees at any point in
    the forest is given by \[\rho(\textrm{distance in miles from that
      point to the stream}).\] The distance from the $k$-th slice to the
    stream is approximately $x_k - (-3) = x_k + 3$, so the density of
    trees in this slice is approximately $\rho(x_k + 3)$.

    We also need to approximate the area of this slice, which we'll do by
    pretending that the slice is a rectangle.  To figure out its height,
    let's go back to the equation of the circle, which is $x^2 + y^2 =
    25$.  Solving for $y$, $y = \pm \sqrt{25 - x^2}$.  The slice goes from
    the bottom half of the circle ($y = -\sqrt{25 - x^2}$) to the top half
    ($y = \sqrt{25 - x^2}$).  So,
    \begin{align*}
      \textrm{height of $k$-th slice} & \approx \sqrt{25 - x_k^2} -
      \left(- \sqrt{25 - x_k^2} \right) = 2 \sqrt{25 - x_k^2} \\
      \textrm{area of $k$-th slice} & \approx (\textrm{height of $k$-th
        slice}) (\textrm{width of $k$-th slice}) \approx 2 \sqrt{25 -
        x_k^2} \Delta x \\
      \textrm{number of trees in $k$-th slice} & \approx
      (\textrm{density of trees in $k$-th slice}) (\textrm{area of
        $k$-th slice}) \\
      & \approx 2 \sqrt{25 - x_k^2} \rho(x_k + 3) \Delta x 
    \end{align*}
    Summing over all the slices and taking the limit as the number of
    slices tends to $\infty$, the number of trees to the right of the
    stream is given by the integral $\boxed{\int_{-3}^5 2 \sqrt{25
        - x^2} \rho(x + 3)~dx}$.
  \end{solution}
  
\end{enumerate}


\end{document}











